% !TEX program = xelatex
\documentclass{英语试卷}
\begin{document}
\examtitle{2024/ 2025 学年佛山市普通高中教学质量检测 (二)}{英语}
\begin{notice}
\par 1. 答卷前, 考生务必用黑色字迹的钢笔或签字笔将自己的姓名、考生号、试室号和座位号填写在答题卡上。用 2B 铅笔将试卷类型 (A) 填涂在答题卡相应位置上, 并在答题卡相应位置上填涂考生号。因笔试不考听力, 试卷从第二部分开始, 试题序号从“21"开始。\par
\par 2. 作答选择题时，选出每小题答案后，用2B铅笔把答题卡对应题目选项的答案信息点除黑；如需改动，用橡皮擦干净后，再选涂其他答案。答案不能答在试卷上。\par
\par 3. 非选择题必须用黑色字迹的钢笔或签字笔作答, 答案必须写在答题卡各题目指定区域内相应位置上,如需改动, 先划掉原来的答案, 然后再写上新答案; 不准使用铅笔和涂改液。不按以上要求作答无效。\par
\par 4. 考生必须保持答题卡的整洁。考试结束后，将试卷和答题卡一并交回。\par
\end{notice}
\parttitle{第二部分 阅读理解（共两节，满分50分）}
\sectitle{第一节（共15小题；每小题2.5分，满分37.5分）}
\par 阅读下列短文，从每题所给的A、B、C和D四个选项中，选出最佳答案。\par
\passagelabel{A}
\par Frequently Asked Questions about London Theatres\par
\par London has a world-leading theatre scene, and it's particularly strong when it comes to family shows. Read the FAQ to make your family theatre trip a pleasant one.\par
\par Can children go to all theatres?\par
\par Most shows are suitable for people aged 12 or older. Some productions might have material that is too mature for younger theatregoers. However, there are also fantastic family shows which are appropriate for all ages, or which are specifically aimed at kids.\par
\par Are children booster seats available?\par
\par Some theatres, though not all, provide booster seats—for example, Nederlander theatres rent height-adjustable seats for £1 if you ask a member of staff. They are available on a first-come, first-served basis. Some venues will also have height restrictions so that audience sitting behind you can still see.\par
\par Are there frequent breaks throughout the duration of a performance?\par
\par If it's a show for a general audience, there will usually be one or two intervals based on the running time. Most shows designed specifically for younger audiences will have short running times, like an hour or less.\par
\par Can I bring my own snacks and drinks into the theatre?\par
\par Most theatres won't allow you to bring your own snacks and drinks. They sell food and drink in the theatre and that's part of how they bring in profit. If you have concerns about dietary requirements, you can contact the theatre beforehand to make arrangements.\par
\par Find all the latest recommendations and more information about the rules on Telegraph Tickets, or check with the box office.\par
\question{21. Where is the text probably from?}
\begin{choices}(2)
  \item A theatre guide.
  \item A drama review.
  \item A show advertisement.
  \item A performance schedule.
\end{choices}
\question{22. Which best describes booster seats?}
\begin{choices}(2)
  \item Found in every theatre.
  \item Provided free of charge.
  \item Required by safety rules.
  \item Offered in limited number.
\end{choices}
\question{23. What are the theatres' concerns behind their food policies?}
\begin{choices}(2)
  \item Financial gain.
  \item Food safety.
  \item Cleaning cost.
  \item Audience experience.
\end{choices}
\passagelabel{B}
\par If you are an enthusiastic reader, you may not have enough room to store all the books you read. And the cost of buying books keeps rising. Now there is a perfect solution in the UK: Bookloop. It is an online buy-back project in partnership with existing booksellers, where readers can trade books they have already read for credit or more pre-owned books.\par
\par Readers can scan (扫描) and upload their books to an online valuation system, and then drop them off at a collection point or have them picked up from their own home, in exchange for credit to be used on other books. Readers must sell at least £5 worth of books to take part. Eric Gagnaire, co-founder of Bookloop, said, "By giving books a second life, we are not only reducing waste but also building a community of readers who value sustainability and support bookshops."\par
\par Bookbag, a bookseller that has joined the program, describes it as "a social space" that gathers a global range of fiction and non-fiction writing. "We're glad to see this program, which feels well-thought-through and a great way of repurposing books and helping sustain bookshops and new book purchase," they said. Authors will also benefit from the program, as royalties from these secondhand book sales will be redistributed through a shared author fund.\par
\par Reception to the launch of Bookloop hasn't been universally positive, however. Some secondhand booksellers like Patrick Kelly, owner of Bookmongers in Brixton, raise concerns about how many books will actually be sold on to readers. Kelly argued that most secondhand books should be recycled. "They are either no longer relevant, overly produced, or are left in bad condition," he said.\par
\par In response, Bookloop stated that only 2\% of collected books end up not being resold-most of the unsold books are donated to local charities, with less than 0.5\% recycled. "Getting pre-owned books to new readers while benefiting both bookshops and authors feels extremely positive," said Gagnaire. "It's a really sustainable way of getting books into the hands of more people, without cost being a barrier."\par
\question{24. What is a requirement for readers to participate in Bookloop?}
\begin{choices}(2)
  \item Trading their credit for new books.
  \item Selling used books worth £5 or more.
  \item Dropping off their books at a local library.
  \item Uploading personal information to the system.
\end{choices}
\question{25. What does Bookbag think of Bookloop?}
\begin{choices}(2)
  \item It is a win-win initiative.
  \item It boosts interest in reading.
  \item It is a well-planned book fair.
  \item It encourages authors to write.
\end{choices}
\question{26. Why is Patrick Kelly mentioned?}
\begin{choices}(2)
  \item To clarify a misunderstanding.
  \item To suggest a brand-new resolution.
  \item To present a contrasting viewpoint.
  \item To highlight an unusual phenomenon.
\end{choices}
\question{27. What can be the best title for the text?}
\begin{choices}(2)
  \item The Rise of Secondhand Book Trading
  \item Reducing Waste, Building a Reputation
  \item Bookloop: A Second Life for Used Books
  \item Innovating Reading: The Bookloop Initiative
\end{choices}
\passagelabel{C}
\par Wildfires are getting worse across the world. Twenty-one years of satellite data shows severe fires have more than doubled in frequency since 2003 and the six most extreme years since records began have come since 2017, causing more deaths and destruction.\par
\par One European Union (EU) pilot program, the Life Landscape Fire Project (LLFP), is now using hungry mountain goats to prevent wildfires before they start. It involves a combination of controlled fires and hungry goats grazing (吃草) to reduce plant cover that can dry up in summer and fuel the rapid spread of fires. Controlled fires help burn older plants that goats don't eat, allowing fresh green vegetation to grow which grazers like goats and cows do eat. The experimental project pays shepherds (牧羊人) to guide their goats to areas where they want the plant cover reduced, and they also build shelters and watering areas to attract free grazers to the desired spots. That's how mountain goats and cows become mobile fire prevention units.\par
\par LLFP director Andre Mota oversees this EU-funded project. "We are in one of the most affected regions by wildfires, historically speaking," said Mota. "Last September we had a wildfire with extreme behavior that burned almost everything in its path. But our parcels managed to keep green." A video from Mota showed that one half of the mountain was black and burned and the other green, the parts which controlled fires and hungry goats had made unlikely to burn. "We believe this proves that, when scaled up from a pilot project to a full solution, it works," he added.\par
\par The sound of goat bells rings out as the goats move peacefully on their way, offering hope in a region directly in the path of climate chaos.\par
\question{28. What do we know about the LLFP?}
\begin{choices}(2)
  \item It protects mountain goats.
  \item It employs integrated methods.
  \item It has been conducted across Europe.
  \item It wins shepherds' unconditional support.
\end{choices}
\question{29. What does the underlined words "the desired spots" refer to?}
\begin{choices}(2)
  \item Free grazers' favorite destinations.
  \item Places far from shepherds' activities.
  \item Spots with abundant water resources.
  \item Areas in need of vegetation reduction.
\end{choices}
\question{30. What is paragraph 3 mainly about?}
\begin{choices}(2)
  \item The influence of wildfires.
  \item The effectiveness of the LLFP.
  \item The challenges faced by the LLFP.
  \item The importance of preventing fires.
\end{choices}
\question{31. What message is conveyed in the last paragraph?}
\begin{choices}(2)
  \item Goats stand as a symbol of peace.
  \item Goats ring the alarm for quick action.
  \item Goats provide hope for wildfire prevention.
  \item Goats offer an ideal solution to climate change.
\end{choices}
\passagelabel{D}
\par How individuals describe their experiences plays a crucial role in shaping emotions and decision-making. However, little research has examined whether language can impact how people experience time alone. To find it out, researchers studied how different terms used to describe time alone affect individuals' understanding of it and their emotions and behaviors.\par
\par In the first study, 500 adults were randomly assigned to evaluate one of five labels describing time alone: metime, time alone, solitude (孤单), being alone, or isolation (孤独). They rated their assigned term on several dimensions, including how positive or negative it felt, its perceived impact on well-being, and whether they actively sought or avoided that type of experience.\par
\par This study revealed that "me-time" was consistently rated as the most positive, with participants considering it as beneficial to well-being and something they actively sought out. In contrast, "isolation" was viewed as the least desirable. The other three labels fell somewhere in between.\par
\par In the second study, 176 undergraduate students were randomly assigned to a 30-minute solitude period described as either me-time or isolation. Before the session, they completed a survey measuring baseline beliefs about being alone. During the session, they avoided all social interactions, but could engage in non-social activities like reading or writing. Afterward, they rated their emotions, described their thoughts and behaviors, and reassessed their beliefs about being alone.\par
\par Those who spent 30 minutes alone after being told they were having "me-time" reported an increase in positive effect, while those who were told they were in "isolation" experienced a decline in positive emotions. Additionally, participants in the "me-time" condition reported more positive beliefs about being alone after the session compared to their baseline beliefs, while those in the "isolation" condition showed no significant change.\par
\par Interestingly, choice of terms did not significantly influence the activities participants engaged in. But it did appear to shape thought content, with those in the "isolation" condition being more likely to think about schoolwork, while those in the "me-time" condition were more likely to reflect on personal growth.\par
\par The studies, however, involved individuals that consisted primarily of adults and undergraduate students, limiting the generalizability of the findings across diverse populations.\par
\question{32. One of the research purposes is to explore \blanknum{}.}
\begin{choices}(2)
  \item the benefits of positive emotions
  \item people's preferences for time alone
  \item how people make decisions while alone
  \item how language shapes people's perceptions
\end{choices}
\question{33. What can be concluded from the first study?}
\begin{choices}(2)
  \item People treasure their time of being alone.
  \item Label selection reflects people's personality.
  \item People prefer being alone to social interaction.
  \item Terms for time alone affect emotional responses.
\end{choices}
\question{34. Why were the baseline beliefs about being alone rated?}
\begin{choices}(2)
  \item To predict future outcomes.
  \item To ensure accurate understanding.
  \item To track potential opinion shifts.
  \item To group the participants accordingly.
\end{choices}
\question{35. What is the limitation of the research?}
\begin{choices}(2)
  \item Contradictory results.
  \item Restricted participant range.
  \item Controlled research duration.
  \item Unreliable research methods.
\end{choices}
\sectitle{第二节（共5小题；每小题2.5分，满分12.5分）}
\par 根据短文内容，从短文后的选项中选出能填入空白处的最佳选项。选项中有两项为多余选项。\par
\par In a conversation, it's often not just what we say but what we're asked that makes the difference. The right question can unlock thoughts and feelings we didn't even know we wanted to share.\par
\par Imagine you're at a party making small talk. Instead of the usual "How's work?" someone asks, "What's been inspiring you lately?" \blanknum{36} You find yourself sharing thoughts about a book that changed your views or a recent experience that moved you. The interaction transforms from polite little chat to genuine engagement. \blanknum{37} In professional settings, thoughtful questions can lead to better problem-solving and stronger teams. During brainstorming, rather than asking "Any ideas?" try "What would we do if resources were unlimited?" Thoughtful questions like these spark creativity and innovation.\par
\par What makes a question "good"? The best questions come from pure curiosity and a desire to understand, not to confirm existing beliefs or direct conversations toward predetermined conclusions.\par
\par \blanknum{38} Before asking something, consider what you really want to learn about the person or topic. It's the difference between asking "How's the new job?" out of politeness and "What's been most surprising about your new role?" from real interest.\par
\par Open-ended questions invite richer responses. Instead of "Did you enjoy your vacation?" ask "What moment from your trip will you remember most vividly?" \blanknum{39}\par
\par Context matters too. Questions appropriate for close friends might be too personal for new acquaintances. Being mindful of relationships and situations helps create comfortable and meaningful exchanges.\par
\par The art of asking questions and establishing deeper connections is an ongoing journey. \blanknum{40} Each one is an invitation to understand someone better, to see the world through their eyes, and to create moments of close human connection. Your next question might unlock an in-depth conversation, a new understanding, or even a life-changing relationship. 微信公众号 IAI English\par
\begin{choices}(1)
  \item Suddenly, the conversation shifts.
  \item You might not know how to respond.
  \item This power extends beyond social situations.
  \item Creating powerful questions starts with genuine curiosity.
  \item A good question makes the other person feel positive and uplifted.
  \item This way, you're encouraging specific experiences and emotions sharing.
  \item As you move forward, challenge yourself to ask more meaningful questions.
\end{choices}
\parttitle{第三部分 语言知识运用（共两节，满分30分）}
\sectitle{第一节（共15小题；每小题1分，满分15分）}
\par 阅读下面短文，从短文后各题所给的A、B、C和D四个选项中，选出可以填入空白处的最佳选项。\par
\par A 12-year-old girl's playground dream has come true, thanks to the kindness of Clearfield City officials who supported her creativity and passion for inclusive play.\par
\par Several months ago, young Rosili noticed that a new playground at Thornock Park of the city was being built, but the \blanknum{41} seemed more suitable for older kids. \blanknum{42} that she wouldn't be able to enjoy it with her younger sisters, Rosili decided to \blanknum{43}.\par
\par Armed with creativity and a passion for play, she mapped out a(n) \blanknum{} 44\blanknum{} playground, complete with detailed descriptions and labels.\par
\par With support from her mom, Rosili took her \blanknum{45} to the city hall, where she met with Eric Howes, the city's community services director. Impressed by her drawings, Howes took time to \blanknum{46} each of them. "She put a lot of thought into that, and she \blanknum{47} it from four different angles," Howes explained.\par
\par Although the Thornock Park playground was already \blanknum{48}, the city had plans to build a new playground at Bicentennial Park, and Rosili's designs were included in the \blanknum{49}.\par
\par Rosili \blanknum{50} every step of the process, working alongside city officials and playground equipment sellers to realize her vision. "Every day I would have to go there for another \blanknum{51}, and I would tell my friends and be like, 'Ah, I get to go again, woooo!'" she \blanknum{52} in social media.\par
\par Finally, Rosili's dream playground was officially opened. Rosili, alongside her family, \blanknum{53} cut the ribbon (彩带) at the ceremony, creating a new play space for children of all ages to enjoy-a space that was \blanknum{54} but brought to life by a young girl who believed she could \blanknum{55}.\par
\cz{41}{park}{city}{community}{equipment}
\cz{42}{Concerned}{Warned}{Surprised}{Confused}
\cz{43}{back out}{step in}{take charge}{make a complaint}
\cz{44}{lively}{traditional}{local}{ideal}
\cz{45}{photos}{designs}{notes}{models}
\cz{46}{praise}{post}{review}{rate}
\cz{47}{recorded}{understood}{showed}{simplified}
\cz{48}{in progress}{out of service}{on display}{at capacity}
\cz{49}{contract}{report}{project}{research}
\cz{50}{took over}{participated in}{worked out}{left behind}
\cz{51}{meeting}{deal}{game}{challenge}
\cz{52}{predicted}{protested}{advertised}{shared}
\cz{53}{proudly}{calmly}{gratefully}{skillfully}
\cz{54}{reserved}{imagined}{occupied}{misused}
\cz{55}{have a dream}{make a choice}{make a difference}{take a risk}
\sectitle{第二节（共10小题；每小题1.5分，满分15分）}
\par 阅读下面短文，在空白处填入1个适当的单词或括号内单词的正确形式。\par
\par Macau draws tourists from every corner of the world with its excellent food, outstanding architecture, history and dynamic culture. This small city in southeast China, recognized as a "Creative City of Gastronomy" by UNESCO, \blanknum{} (offer) a lot for visitors to experience.\par
\par If you are unsure about what sightseeing spots to explore, start with the St. Paul Ruins and Senado Square, \blanknum{57} are the most famous landmarks and the most photographed places in the city and must-sees for those \blanknum{58} (visit) Macau for the first time. From there, follow your inner sense and be adventurous, and Macau will reward you \blanknum{59} its true beauty and hidden spots. And don't hesitate to taste the local street food you will find on the way.\par
\par Macau's integration of diverse cultures \blanknum{60} (reflect) in its numerous churches and temples, which, like most of the landmarks in town, don't charge for the \blanknum{61} (enter). So, feel free to walk in and discover.\par
\par No trip to Macau would be complete without taking \blanknum{62} tour of the historical Taipa. The picturesque \blanknum{63} (street) of Taipa Village are filled with \blanknum{64} (vivid) colored houses, stunning churches, ancient temples and traditional Chinese shops. \blanknum{65} (experience) a taste of Macanese bites and treats, just eat and drink your way around. You'll never get stuck for ideas on where to dine at Taipa Village!\par
\parttitle{第四部分 写作（共两节，满分40分）}
\sectitle{第一节（满分15分）}
\par 你校组织全体学生每天上午大课间集体跑步，你认为这样的安排不合理。请你写一篇短文向校英语报"Your Say"栏目投稿，内容包括：\par
\par (1) 说明你的顾虑;\par
\par (2) 提出建议。\par
\par (1) 写作词数应为 80 左右;\par
\par (2) 请按如下格式在答题卡的相应位置作答。\par
\par Time to Rethink Our Long Break Routine\par
\sectitle{第二节（满分25分）}
\par 阅读下面材料，根据其内容和所给段落开头语续写两段，使之构成一篇完整的短文。\par
\par When Jane taught the sixth-grade class at Morris School, all 34 of her students were dear to her, while Mark was truly special. He had that happy-to-be-alive attitude that made even his occasional tricks delightful.\par
\par Mark talked all the time. Jane reminded him again and again that talking without permission was unacceptable, and his response always sounded sincere: "Thank you for correcting me, Sister!"\par
\par One morning, Jane's patience was growing thin when Mark talked once too often. She looked at Mark and said, "If you say one more word, I'm going to tape your mouth shut!" Not even ten seconds passed before someone said, "Mark is talking again." Having stated the punishment in front of the class, Jane decided she had to act on it. She walked to Mark's desk, tore off two pieces of tape, made a big X with them over his mouth, and then returned to the front of the classroom.\par
\par As Jane glanced at Mark to see how he was doing, he winked (眨眼) at her. That did it! Jane started laughing. The entire class cheered as she shook her head and walked back to remove the tape. His first words were, "Thank you for correcting me, Sister."\par
\par One Friday afternoon, after working hard on a new idea all week, the students were getting discouraged and starting to get annoyed with each other. To help lift their spirits and encourage a positive atmosphere, Jane asked them to write down the nicest things about the other students on a piece of paper and hand them in. Mark appeared nervous as if he was worrying his classmates would forget anything nice about him. That Saturday, Jane wrote down each student's name on a separate sheet of paper and listed what everyone else had said about them.\par
\par On Monday, Jane gave each student their lists. As the children read them, the classroom was filled with quiet excitement and soft laughter.\par
\par 注意：\par
\par （1）续写词数应为150左右；\par
\par （2）请按如下格式在答题卡的相应位置作答。\par
\writingpara{Paragraph 1: Mark's eyes widened as he read his paper.}
\writinglines{8}
\writingpara{Paragraph 2: Jane was overjoyed to see Mark again at school after so many years.}
\writinglines{8}
\end{document}
